\chapter{Discussion}
\startcontents[chapters]

One of the surprising outcomes of the study is that there is no monotonous change from one value of gender to the other when going from attributive to relative pronoun agreement, which is contrary to one of the hypotheses. This may be a result of using low quality data for analysis -- the corpus, from which agreement was extracted, was first parsed automatically and the dependency trees produced this way have not been reviewed. It is possible that some of the words are annotated incorrectly or that there are non-existent dependencies found in the sentences.

It is also worth pointing out that there were very few instances of relative pronoun agreement -- 208 compared to 1478 and 1873 instances of predicate and attributive agreement respectively. In the univariable model predicting by type of agreement the difference between the relative pronoun and the predicate agreement is less significant than the difference between the predicate and attributive agreement, while the relative pronoun predictor in the multivariable model is not significant altogether. The unexpected decline in the proportion of semantic agreement between predicate and relative pronouns can be a result of limited relative pronoun data considered in the analysis. One of the possible directions for future research is to focus on hybrid agreement with relative and personal pronouns, since the current study offered little insight into those agreement targets.

There is also a potentially interesting observation to be made about the large mean distance between the target and controller of agreement, specifically in predicate agreement. The distance was found to be a significant predictor for semantic agreement (the larger the distance the more likely it is that semantic agreement will be chosen), but only by the univariable model. The multivariable model shows, that the effect of distance stops being statistically significant when other predictors are taken into account: one of the significant ones is predicative agreement. As a future perspective, it would be beneficial to test the predicative agreement with small distances between targets and controllers, to see if the preference for semantic agreement was here caused by large distance or by predicate targets.

\cite{corbett_agreement_1979} observed the ordering of the target and controller of agreement to have an influence over the choice between semantic and syntactic agreement -- this was found in the current study as well. According to both models, univariable and multivariable, if the target of agreement is placed before the controller it is unlikely that the agreement will be semantic. The effect is highly significant for both models.

Finally, another hypothesis tested here was that split hybrids are less likely to have semantic agreement than lexical hybrids. Depreciatives, which are an example of split hybrids, had semantic agreement in over 20\% of all occurrences, while the socially gendered nouns had semantic agreement in only 2.85\%, which is contrary to the expectations. It could be interesting to compare the depreciative forms with different split hybrids, to see it this preference for semantic agreement is exclusive to depreciatives or not. In an experimental study it would be possible to test the agreement with the Polish \textsl{ręka} (meaning 'hand') used by \cite{corbett_agreement_2023} to introduce split hybrids, since the examples in the corpus were not sufficient to do so in the present study. 

The corpus analysis yielded many occurrences of agreements with profession nouns, but unfortunately many of those were not useful for the present study. Without manually selecting the occurrences it was impossible to differentiate the examples where the referent was male from the ones where the referent was female, but all agreements were masculine. In the latter case it would have been an instance of hybrid agreement, where the speaker simply chose to use syntactic agreement for every target -- an interesting example that is technically possible, but could not have been selected automatically from the whole database of profession noun agreement. Nevertheless, the socially gendered nouns provided the opportunity to study Polish lexical hybrids, so the study was not harmed.

What should be noted, is that 15 of the profession nouns did not have any feminine agreements at all. Most of those words appeared in the corpus less than a 100 times each, the median of their number of occurrences in the corpus was 37. Since those nouns are rare, a better option would be to study them experimentally. There are no obvious feminitives for them \citep{szpyra-kozlowska_premiera_2019, szpyra-kozlowska_czy_2022}, therefore there is a high chance of finding hybrid agreement with them.
